\chapter{Parameter methods: Feynman, Frullani, Leibniz}

\begin{orient}
\textbf{What this chapter is for.} The most powerful move in definite integration is to make the problem harder on purpose. Given an integral with no free constant in it, you invent one, producing a whole family of integrals of which the original is a single member. You then differentiate with respect to the invented constant, which turns the integrand into something you can actually do, integrate the result back up, and pin the constant of integration by evaluating the family at a value of the parameter where the answer is obvious. The chapter treats this as a six step protocol and insists on the two steps every student drops: proving that the derivative may be taken inside the integral, and choosing an anchor value at which the family is genuinely known. Around that core sit its relatives: the Leibniz rule when the limits themselves move, Frullani's formula for a difference of two scalings, parameter differentiation of the Beta function, and the reverse move of replacing a blocking factor by an integral representation. At the end you should be able to look at an integral with no parameter and see where to put one.
\end{orient}

\section{Making the problem harder on purpose}

An indefinite integral is a search for an antiderivative. A definite integral is not: it is a single number, and the number can often be reached without ever finding an antiderivative. Parameter methods are the largest family of such routes. They work because differentiation with respect to a parameter is a different operation from differentiation with respect to $x$, and it acts on the integrand in a way that can destroy exactly the feature that made the integral impossible. A $\ln x$ in a denominator is immovable under any substitution; it is produced for free by $\partial_a x^{a}$. A factor $1/x$ that ruins convergence at the origin is produced for free by $\partial_a$ acting on an exponent. A $\ln\sin x$ that no trigonometric identity will absorb is produced for free by $\partial_p$ acting on $\sin^{2p-1}x$.

The price is that you must solve a family of problems rather than one. That is almost always a bargain, because the family carries information that the single member does not. A family satisfies a differential equation in $a$, and a differential equation plus one boundary value determines everything. The single member satisfies nothing.

The names attached to these methods are historical accidents worth knowing so that the literature is readable. Differentiation under the integral sign is Leibniz's rule, and the version with moving limits is the one Leibniz actually proved; it acquired Feynman's name from a single paragraph in his memoirs describing how a book of tricks he read as a boy let him evaluate integrals nobody else in the room could. Frullani's integral dates from 1821 and was made rigorous by Hardy a century later, which is why its hypotheses come in two versions. Schwinger's parametrisation is the youngest of the group and comes from quantum field theory, where products of many denominators are the normal case. None of the attributions matters mathematically, but all four names appear in problem statements.

There is also a psychological obstacle worth naming. Generalising a problem feels like retreat, and students trained on antiderivatives will spend an hour on substitutions before they will spend a minute inventing a parameter. The correct instinct is the opposite: as soon as an integral has a definite value but no visible antiderivative, ask what family it belongs to. Almost every named definite integral in the subject is the special value of a family, and the family is usually easier than the value.

Here is the protocol in full. Steps 4 and 6 are the ones that get skipped, and they are the ones that produce wrong answers.

\begin{keynote}[The six steps]
\textbf{1.} Introduce a parameter $a$, so that $I=I(a_\star)$ for a specific $a_\star$, and $I(a)$ is a family.
\textbf{2.} Differentiate: $I'(a)=\int \pdv{f}{a}\dd x$. If this is not easier than $I$, you inserted the parameter in the wrong place; go back to step 1.
\textbf{3.} Evaluate $I'(a)$ in closed form.
\textbf{4.} Justify the interchange, on a closed parameter interval containing $a_\star$: exhibit $g$ with $\abs{\pdv{f}{a}}\leq g(x)$ for all $a$ in that interval and $\int g<\infty$.
\textbf{5.} Integrate back: $I(a)=\int I'(a)\dd a+C$.
\textbf{6.} Fix $C$ from an anchor $a_0$ at which $I(a_0)$ is known, and check that $a_0$ lies in the region where steps 2 to 4 were valid.
\end{keynote}

The first technique is the protocol itself, executed on an integral that is genuinely inaccessible without it.

\tech{Differentiation under the integral sign}{core}
\cue A constant in the integrand that could obviously have been a variable: $\ln(a+x)$, $\eu^{-ax}$, $x^{a}$, $\arctan(ax)$, $(1+a^{2}x^{2})$. Equivalently, a specific number in the problem which, if replaced by a letter, would make the integral trivial at some other value of that letter.
\method Set $I(a)=\int_c^d f(x,a)\dd x$, compute
\[I'(a)=\int_c^d\pdv{f}{a}\dd x,\]
evaluate it, integrate back in $a$, and fix the constant at an anchor. The whole method rests on $\partial_a$ producing an integrand that is elementary even though $f$ was not.

\begin{worked}
\[I(a)=\int_0^1\ln\!\left(\frac{a+x}{a-x}\right)\frac{\dd x}{x\sqrt{1-x^{2}}},\qquad a>1.\]
Nothing here is an antiderivative problem. Differentiate instead. Since $\partial_a\bigl[\ln(a+x)-\ln(a-x)\bigr]=\tfrac{1}{a+x}-\tfrac{1}{a-x}=\tfrac{-2x}{a^{2}-x^{2}}$, the awkward $1/x$ cancels:
\[I'(a)=-2\int_0^1\frac{\dd x}{(a^{2}-x^{2})\sqrt{1-x^{2}}}=-2\int_0^{\pi/2}\frac{\dd\theta}{a^{2}-\sin^{2}\theta},\]
using $x=\sin\theta$. Divide numerator and denominator by $\cos^{2}\theta$ and put $t=\tan\theta$: the standard value is $\int_0^{\pi/2}\frac{\dd\theta}{p+q\sin^{2}\theta}=\frac{\pi}{2\sqrt{p(p+q)}}$, so with $p=a^{2}$, $q=-1$,
\[I'(a)=-\frac{\pi}{a\sqrt{a^{2}-1}}.\]
Now integrate back. On $a>1$, $\dvop{a}\arccos(1/a)=\frac{1}{a\sqrt{a^{2}-1}}$, so $I(a)=-\pi\arccos(1/a)+C$. The anchor is $a\to\infty$: the integrand tends to $0$ pointwise and is dominated by an integrable function, so $I(\infty)=0$, and since $\arccos 0=\tfrac{\pi}{2}$ we get $C=\tfrac{\pi^{2}}{2}$. Hence
\[I(a)=\frac{\pi^{2}}{2}-\pi\arccos\frac1a,\qquad\text{and at }a=2,\quad I(2)=\frac{\pi^{2}}{2}-\frac{\pi^{2}}{3}=\frac{\pi^{2}}{6}\approx1.644934.\]
The licence for step 2 is visible: for $a\geq a_0>1$, $\bigl\lvert\partial_a f\bigr\rvert=\frac{2}{(a^{2}-x^{2})\sqrt{1-x^{2}}}\leq\frac{2}{(a_0^{2}-1)\sqrt{1-x^{2}}}$, and that dominating function has finite integral $\frac{\pi}{a_0^{2}-1}$ on $[0,1]$. Checking the general formula at a second point, $a=\tfrac32$ gives $\tfrac{\pi^{2}}{2}-\pi\arccos\tfrac23\approx2.292507$, which quadrature confirms.
\end{worked}

\begin{trap}
Two failures, both at step 6. The first is forgetting $C$ altogether and reporting $-\pi\arccos(1/a)$, which is off by $\pi^{2}/2$ and negative. The second is anchoring at a value where the family does not exist: in the example above $a=1$ looks tempting because $\arccos 1=0$, but the integral diverges there and the derivative formula breaks down, so no information can be extracted from it.
\end{trap}

\begin{seealsob}Choosing where to insert the parameter; the parameter ODE; Beta and Gamma parameter differentiation.\end{seealsob}

The example was chosen so that the parameter was already present. The real skill is the opposite case, where the integral has no constant in it at all and you must manufacture one. This is where the method becomes a genuine art, because a bad insertion makes $I'(a)$ harder than $I$ and there is no warning bell.

\tech{Choosing where to insert the parameter}{hard}
\cue A hard integral with no free constant: $\int_0^1\frac{x-1}{\ln x}\dd x$, $\int_0^\infty\frac{\sin x}{x}\dd x$, $\int_0^1\frac{\arctan x}{x}\dd x$. There is nothing to differentiate with respect to until you put something there.
\method Insert the parameter where $\partial_a$ \emph{kills the obstruction}. Three insertions cover most cases. Replace an exponent by $a$: $\partial_a x^{a}=x^{a}\ln x$, which cancels a $\ln x$ downstairs. Attach a damping factor $\eu^{-ax}$: $\partial_a$ brings down $-x$, which cancels a $1/x$ and simultaneously repairs convergence at infinity. Generalise a coefficient already in the problem, as $\arctan x\mapsto\arctan(ax)$. Then check step 2 before committing: if $\partial_a f$ is not visibly easier, choose again.

\begin{worked}
\[\int_0^1\frac{x-1}{\ln x}\dd x.\]
The obstruction is $\ln x$ in the denominator, so the insertion goes in the exponent: put $I(a)=\int_0^1\frac{x^{a}-1}{\ln x}\dd x$, which recovers the target at $a=1$. Then
\[I'(a)=\int_0^1\frac{x^{a}\ln x}{\ln x}\dd x=\int_0^1x^{a}\dd x=\frac{1}{a+1},\]
so $I(a)=\ln(a+1)+C$. The anchor is $a=0$, where the numerator is identically zero and hence $I(0)=0$, forcing $C=0$. Therefore $I(a)=\ln(a+1)$ and the answer is $\ln2\approx0.693147$. The same computation with two parameters gives $\int_0^1\frac{x^{a}-x^{b}}{\ln x}\dd x=\ln\frac{a+1}{b+1}$; at $a=3$, $b=1$ this predicts $\ln2$ again, which quadrature confirms.

Contrast a bad insertion on the same problem. Writing $J(a)=\int_0^1\frac{x-1}{\ln(ax)}\dd x$ also recovers the target at $a=1$, but $\partial_a$ produces $-\frac{x-1}{a\ln^{2}(ax)}$, which has a worse singularity than the original and no closed form. The parameter went into the right function and the wrong slot.

A third possibility is worse than either, because it looks like progress. For $\int_0^1\frac{\arctan x}{x}\dd x$ the natural insertion is $K(a)=\int_0^1\frac{\arctan(ax)}{x}\dd x$, and indeed $K'(a)=\int_0^1\frac{\dd x}{1+a^{2}x^{2}}=\frac{\arctan a}{a}$, which is elementary. But integrating back gives $K(1)=\int_0^1\frac{\arctan a}{a}\dd a$, which is the original integral again. The parameter was inserted along the direction the integral is invariant under, so the protocol returns an identity rather than a value. The correct diagnosis is that this integral is not a parameter problem at all: its value is Catalan's constant $G$, reached by series expansion in the next chapter.
\end{worked}

\begin{center}
\begin{tikzpicture}[node distance=6mm and 8mm]
\node[dtest] (a) {Is a $\ln x$ or $\ln$ factor\\ blocking the integral?};
\node[dleaf, below left=9mm and 4mm of a] (b) {Put $a$ in an exponent:\\ $x\mapsto x^{a}$, since\\ $\partial_a x^{a}=x^{a}\ln x$};
\node[dtest, below right=9mm and 4mm of a] (c) {Does the tail fail to\\ converge absolutely?};
\node[dleaf, below left=9mm and 0mm of c] (d) {Attach $\eu^{-ax}$;\\ $\partial_a$ supplies $-x$\\ and fixes the tail};
\node[dtest, below right=9mm and 0mm of c] (e) {Is there a constant\\ already in the problem?};
\node[dleaf, below left=9mm and 0mm of e] (f) {Vary it, then verify\\ $\partial_a f$ is easier};
\node[dleaf, below right=9mm and 0mm of e] (g) {Replace a factor by an\\ integral representation};
\draw[dedge] (a) -- node[dlab,left]{yes} (b);
\draw[dedge] (a) -- node[dlab,right]{no} (c);
\draw[dedge] (c) -- node[dlab,left]{yes} (d);
\draw[dedge] (c) -- node[dlab,right]{no} (e);
\draw[dedge] (e) -- node[dlab,left]{yes} (f);
\draw[dedge] (e) -- node[dlab,right]{no} (g);
\end{tikzpicture}
\end{center}

The tree routes on the obstruction, not on the integrand, and that is the point. Two integrals that look nothing alike take the same branch if the same feature is blocking them, and the same integral takes different branches depending on which of its features you decide is the obstacle. Read the leftmost branch first: a logarithm in a denominator is the most common blocker and the exponent insertion is the most reliable cure. The rightmost leaf is a genuine fallback rather than a technique of last resort; it appears again as the first technique of the next chapter, where the added integral is turned into an added sum.

\begin{trap}
Inserting the parameter so that $I'(a)$ is harder than $I(a)$, and then persevering. There is no rule that a family must be easier than its member. The second trap is assuming the anchor: $I(0)=0$ held above because $x^{0}-1$ vanishes identically, but in $\int_0^\infty \eu^{-ax}\frac{\sin x}{x}\dd x$ the value at $a=0$ is precisely what you are trying to find, and using it as an anchor is circular.
\end{trap}

\begin{seealsob}Differentiation under the integral sign; Frullani's integral; Schwinger parametrisation.\end{seealsob}

\begin{keynote}[Insertions that repay themselves]
Four insertions account for most successful applications, and each is chosen for what $\partial_a$ produces rather than for what looks natural. Exponent, $x\mapsto x^{a}$, produces $\ln x$ and is the answer whenever a logarithm sits in a denominator. Damping, $f\mapsto \eu^{-ax}f$, produces a factor $-x$ and is the answer whenever a $1/x$ blocks convergence or an oscillatory tail refuses to settle. Scale inside a bounded function, $\arctan x\mapsto\arctan(ax)$ or $\ln(1+x)\mapsto\ln(1+ax)$, produces a rational function and is the answer when the obstruction is the transcendental function itself. Quadratic weight, $\eu^{-x^{2}}\mapsto \eu^{-ax^{2}}$, produces $-x^{2}$ and converts a Gaussian family into a power of $a$. If none of the four applies, the integrand probably wants an integral representation instead of a derivative.
\end{keynote}

\begin{keynote}[Rescale before you differentiate]
A parameter that enters only through a scaling can usually be removed before any differentiation, and removing it is free information. If $f(x,a)=g(ax)$ on $(0,\infty)$, then $x\mapsto x/a$ gives $I(a)=\tfrac1a\int_0^\infty g$, so $I$ is determined up to one number and there is nothing to differentiate. The useful cases are the ones where scaling almost works: $\int_0^\infty \eu^{-ax^{2}}\dd x=\tfrac12\sqrt{\pi/a}$ by scaling alone, while $\int_0^\infty \eu^{-ax^{2}}\cos(bx)\dd x$ needs one genuine differentiation because two parameters cannot both be scaled away. Always check for a scaling first: it costs one substitution and can retire the problem.
\end{keynote}

\section{Licensing the interchange, and choosing the anchor}

Steps 4 and 6 are not bureaucracy. They are where parameter methods actually fail, and a reader who treats them as decoration will eventually produce a confident wrong number. Take them in turn.

The interchange $\dvop{a}\int f=\int\partial_a f$ is a limit swap, and limit swaps require a hypothesis. The usable version is this: if $f(x,a)$ and $\partial_a f(x,a)$ are continuous on $[c,d]\times[\alpha,\beta]$, or more generally if there is a single integrable $g$ with $\abs{\partial_a f(x,a)}\leq g(x)$ for every $a\in[\alpha,\beta]$, then $I$ is differentiable on $[\alpha,\beta]$ and the swap is valid there. Two features of that statement are load bearing. The interval in $a$ must be closed, because the theorem is proved by applying dominated convergence to difference quotients and those need a uniform bound. And the dominating function must not depend on $a$: a bound like $\abs{\partial_a f}\leq h(x,a)$ with $\int h(\cdot,a)<\infty$ for each fixed $a$ proves nothing at all.

In practice you produce $g$ by bounding the $a$ dependent parts away from their bad values, which is exactly why the licence is stated on a subinterval $[\alpha,\beta]$ rather than on the whole range of $a$ for which $I(a)$ happens to converge. The subinterval is then enlarged by covering: since differentiability is local, a valid licence on every $[\alpha,\beta]$ inside an open interval gives a valid formula on the open interval. What it never gives you is the endpoint, and the endpoint is usually where the anchor lives.

\begin{keynote}[Where the licence typically fails]
Almost always at an endpoint of the parameter range, and almost always the endpoint you wanted to use as an anchor. In $I(a)=\int_0^\infty \eu^{-ax}\frac{\sin x}{x}\dd x$ the derivative is $I'(a)=-\int_0^\infty \eu^{-ax}\sin x\dd x=-\frac{1}{1+a^{2}}$, and the dominating function $\eu^{-a_0x}$ exists for $a\geq a_0>0$ but not for $a\geq0$: at $a=0$ the differentiated integral $\int_0^\infty\sin x\dd x$ does not converge at all. So the ODE is valid only on $a>0$, giving $I(a)=\tfrac{\pi}{2}-\arctan a$ after anchoring at $I(\infty)=0$, and the value at $a=0$ must be obtained by a \emph{separate} argument: $I$ is continuous at $0$ from the right, because $\int_R^\infty \eu^{-ax}\frac{\sin x}{x}\dd x$ is small uniformly in $a\geq0$ by Dirichlet's test. Only then may you write $\int_0^\infty\frac{\sin x}{x}\dd x=\tfrac{\pi}{2}$.
\end{keynote}

Now the anchor. You need one value $a_0$ at which $I(a_0)$ is known independently, and it must lie in the closure of the region where the derivative computation was licensed. Three anchors cover nearly everything. The first is a value at which the integrand vanishes identically, so that $I(a_0)=0$ for trivial reasons; $a=0$ in $\int_0^1\frac{x^{a}-1}{\ln x}\dd x$ is this case. The second is $a\to\infty$ with an exponential or algebraic decay that lets dominated convergence give $I(\infty)=0$. The third is a value at which the family collapses to a standard integral you already know.

A worked illustration of the second anchor, which also shows a clean partial fraction step. With $I(a)=\int_0^\infty\frac{\ln(1+a^{2}x^{2})}{1+x^{2}}\dd x$ for $a\geq0$,
\[I'(a)=\int_0^\infty\frac{2ax^{2}}{(1+a^{2}x^{2})(1+x^{2})}\dd x=\frac{2a}{a^{2}-1}\int_0^\infty\left[\frac{1}{1+x^{2}}-\frac{1}{1+a^{2}x^{2}}\right]\dd x=\frac{2a}{a^{2}-1}\left(\frac{\pi}{2}-\frac{\pi}{2a}\right)=\frac{\pi}{1+a},\]
where the apparently singular factor $1/(a^{2}-1)$ cancels against $\tfrac{\pi}{2}-\tfrac{\pi}{2a}$. Anchoring at $a=0$, where the integrand is identically zero, gives $I(a)=\pi\ln(1+a)$; at $a=3$ this predicts $\pi\ln4\approx4.355172$, which quadrature returns. The two parameter version costs nothing extra:
\[\int_0^\infty\frac{\ln(1+a^{2}x^{2})}{b^{2}+x^{2}}\dd x=\frac{\pi}{b}\ln(1+ab),\]
checked at $a=2$, $b=3$ against the numerical value $2.037752$.

Not every family has a single formula, and the anchor is what tells you so. Consider $L(a)=\int_0^{\pi}\ln\bigl(1-2a\cos x+a^{2}\bigr)\dd x$. Differentiating and using the substitution $t=\tan\tfrac x2$ gives $L'(a)=0$ for $\abs{a}<1$ and $L'(a)=\tfrac{2\pi}{a}$ for $\abs{a}>1$, so $L$ is constant on each of the two regions separately. The anchor $a=0$ gives $L\equiv0$ on $\abs{a}<1$; the anchor $a\to\infty$, where $\ln(1-2a\cos x+a^{2})=2\ln\abs{a}+O(1/a)$, gives $L(a)=2\pi\ln\abs{a}$ on $\abs{a}>1$. Numerically $L(0.5)=0$ and $L(3)=2\pi\ln3\approx6.902785$, both confirmed. A single anchor would have produced one of the two branches and silently misreported the other.

\begin{keynote}[Uniform convergence or domination]
Two licences are in circulation and they are not the same. Uniform convergence of $\int_c^{d}$ as $d\to\infty$, uniformly in $a$ on a closed interval, is what you use when the trouble is the range of integration: it is proved by a tail estimate and it is the natural tool for oscillatory integrands, where Dirichlet's test supplies the uniformity. A dominating function is what you use when the trouble is the integrand: it is proved by an inequality and it is the natural tool for singularities at a finite endpoint. Neither implies the other. When both are available, prefer domination, because it also gives continuity of $I$ in $a$ for free, and continuity is what lets you push the answer to the endpoint of the parameter range.
\end{keynote}

\subsection{An anchor that is not an endpoint}

The third kind of anchor, a value at which the family collapses to something known, is the one people forget exists, and it is often the cleanest. Consider
\[I(a)=\int_0^{\pi/2}\ln\bigl(a^{2}\cos^{2}x+b^{2}\sin^{2}x\bigr)\dd x,\qquad a,b>0,\]
with $b$ held fixed. Differentiating in $a$ and writing $D=a^{2}\cos^{2}x+b^{2}\sin^{2}x$, the identity $b^{2}-D=(b^{2}-a^{2})\cos^{2}x$ turns $\frac{\cos^{2}x}{D}$ into $\frac{1}{b^{2}-a^{2}}\bigl[\frac{b^{2}}{D}-1\bigr]$, and with the standard value $\int_0^{\pi/2}\frac{\dd x}{D}=\frac{\pi}{2ab}$,
\[I'(a)=2a\int_0^{\pi/2}\frac{\cos^{2}x}{D}\dd x=\frac{2a}{b^{2}-a^{2}}\left(\frac{\pi b}{2a}-\frac{\pi}{2}\right)=\frac{\pi}{a+b}.\]
So $I(a)=\pi\ln(a+b)+C$. Neither $a=0$ nor $a\to\infty$ is usable: at $a=0$ the integrand has a logarithmic singularity at $x=0$ and the derivative formula fails there, and as $a\to\infty$ the integral diverges. But at $a=b$ the family collapses, since $D\equiv b^{2}$ and $I(b)=\pi\ln b$. That forces $C=-\pi\ln2$ and
\[\int_0^{\pi/2}\ln\bigl(a^{2}\cos^{2}x+b^{2}\sin^{2}x\bigr)\dd x=\pi\ln\frac{a+b}{2},\]
checked at $a=2$, $b=5$ against the numerical value $3.935671$. At $a=b=1$ it degenerates correctly to $\int_0^{\pi/2}\ln1\dd x=0$, and setting $a=1$, $b\to0$ recovers $\int_0^{\pi/2}\ln\cos x\dd x=-\tfrac{\pi}{2}\ln2$, a limit that is worth taking as a sanity check on the whole formula.

\subsection{The protocol run end to end}

It is worth watching all six steps once, in order, on a problem where none of them is trivial. Take
\[T=\int_0^{\pi/2}\frac{\arctan(3\tan x)}{\tan x}\dd x.\]

\textbf{Step 1.} The constant $3$ is the obvious letter. Put $T(a)=\int_0^{\pi/2}\frac{\arctan(a\tan x)}{\tan x}\dd x$, so that $T=T(3)$.

\textbf{Step 2.} Differentiating in $a$ multiplies the integrand by $\frac{\tan x}{1+a^{2}\tan^{2}x}$ and cancels the $\tan x$ downstairs, leaving
\[T'(a)=\int_0^{\pi/2}\frac{\dd x}{1+a^{2}\tan^{2}x}=\int_0^{\pi/2}\frac{\cos^{2}x}{\cos^{2}x+a^{2}\sin^{2}x}\dd x.\]
That is elementary where the original was not, so the insertion is the right one.

\textbf{Step 3.} Write $D=\cos^{2}x+a^{2}\sin^{2}x$ and note $a^{2}-D=(a^{2}-1)\cos^{2}x$, so $\frac{\cos^{2}x}{D}=\frac{1}{a^{2}-1}\bigl[\frac{a^{2}}{D}-1\bigr]$. Since $\int_0^{\pi/2}\frac{\dd x}{D}=\frac{\pi}{2a}$,
\[T'(a)=\frac{1}{a^{2}-1}\left(\frac{\pi a}{2}-\frac{\pi}{2}\right)=\frac{\pi}{2(a+1)}.\]

\textbf{Step 4.} On $0\leq a\leq A$ the integrand of $T'$ is bounded by $1$, an integrable function on a set of finite measure, uniformly in $a$. The licence holds on every such interval, hence on $[0,\infty)$.

\textbf{Step 5.} Integrating, $T(a)=\frac{\pi}{2}\ln(1+a)+C$.

\textbf{Step 6.} At $a=0$ the integrand is identically zero, so $T(0)=0$ and $C=0$. Therefore $T(a)=\frac{\pi}{2}\ln(1+a)$ and
\[T=\frac{\pi}{2}\ln4=\pi\ln2\approx2.177586,\]
which quadrature confirms. Notice that step 4 was cheap here precisely because step 3 produced a bounded integrand; whenever the differentiated integrand is bounded on a finite interval, the licence is automatic and should be stated in one line rather than skipped.

\begin{keynote}[Two parameters, and mixed partials]
When one parameter is not enough, use two. The rule is that $\partial_a\partial_b I$ must be easier than either single derivative, and the reconstruction then needs two anchors: integrate in $b$ and fix a function of $a$, then integrate in $a$ and fix a number. Equality of mixed partials gives a free consistency check that is worth performing: compute $\partial_a\partial_b$ and $\partial_b\partial_a$ separately and confirm they agree. If they do not, one of the two interchanges was illegitimate, and you have found the error before it reached the answer.
\end{keynote}

\subsection{A gallery of parameter families}

The following families are worth knowing by sight, because a large fraction of hard definite integrals are members of one of them. Each is proved by the six step protocol, and each has been checked numerically at a specific parameter value.
\[\int_0^{\infty}\eu^{-ax}\frac{\sin(bx)}{x}\dd x=\arctan\frac{b}{a},\qquad \int_0^{\infty}\frac{\eu^{-ax}-\eu^{-bx}}{x}\dd x=\ln\frac{b}{a}\qquad(a,b>0).\]
\[\int_0^{\infty}\frac{\ln(1+a^{2}x^{2})}{b^{2}+x^{2}}\dd x=\frac{\pi}{b}\ln(1+ab),\qquad \int_0^{\infty}\frac{\arctan(ax)}{x(1+x^{2})}\dd x=\frac{\pi}{2}\ln(1+a).\]
\[\int_0^{\pi/2}\frac{\arctan(a\tan x)}{\tan x}\dd x=\frac{\pi}{2}\ln(1+a),\qquad \int_0^{\pi/2}\ln\bigl(a^{2}\cos^{2}x+b^{2}\sin^{2}x\bigr)\dd x=\pi\ln\frac{a+b}{2}.\]
\[\int_0^{\infty}\eu^{-ax^{2}}\cos(bx)\dd x=\frac12\sqrt{\frac{\pi}{a}}\,\eu^{-b^{2}/4a},\qquad \int_0^{\infty}\frac{\sin(ax)}{x(1+x^{2})}\dd x=\frac{\pi}{2}\bigl(1-\eu^{-a}\bigr)\quad(a>0).\]
Three observations are worth extracting. First, $\ln(1+a)$ appears whenever $I'(a)$ is $\frac{c}{1+a}$, which happens whenever the differentiated integrand is a rational function whose partial fractions telescope; this is the single most common shape in the whole subject. Second, several of these families are the same family in disguise: $x=\tan\theta$ carries the second line into the third. Third, the last entry has a restriction $a>0$ that is not cosmetic, since the integrand is odd in $a$, and a formula quoted without its restriction will be used at $a<0$ and will be wrong by a sign.

\section{When the derivative reproduces the integral}

Sometimes $I'(a)$ is not an integral you can evaluate, but integration by parts in $x$ turns it back into $I(a)$ itself. That is not a failure. It is a first order linear differential equation in the parameter, and one boundary value solves it. This is the single most elegant corner of the subject, because the integral is determined by its own structure rather than by any antiderivative, and the calculation never touches the function whose integral you want.

\tech{The parameter ODE}{hard}
\cue $I'(a)$ resists direct evaluation, but the differentiated integrand is the original integrand multiplied by $x$, or by a factor that integration by parts can convert back. Gaussian times trigonometric is the archetype, as are integrands built from $\eu^{-x^{2}}$, $\eu^{-x}x^{a}$, and $\sech x$.
\method Derive a relation $I'(a)=\lambda(a)I(a)+\mu(a)$ by parts in $x$, discarding boundary terms only after checking them. Solve the linear ODE with an integrating factor, and fix the single constant from the undamped or undeformed integral $I(0)$.

\begin{worked}
\[F(t)=\int_0^{\infty}\eu^{-x^{2}}\cos(tx)\dd x.\]
Differentiating in $t$ gives $F'(t)=-\int_0^\infty x\,\eu^{-x^{2}}\sin(tx)\dd x$, licensed because $\abs{x\eu^{-x^{2}}\sin(tx)}\leq x\eu^{-x^{2}}$ uniformly in $t$, and $x\eu^{-x^{2}}$ is integrable on $(0,\infty)$. Integrate by parts with $u=\sin(tx)$ and $\dd v=x\eu^{-x^{2}}\dd x$, so $v=-\tfrac12\eu^{-x^{2}}$:
\[F'(t)=\Bigl[\tfrac12\eu^{-x^{2}}\sin(tx)\Bigr]_0^{\infty}-\frac{t}{2}\int_0^\infty \eu^{-x^{2}}\cos(tx)\dd x=-\frac{t}{2}F(t).\]
The boundary term vanishes at both ends, at infinity because $\eu^{-x^{2}}$ crushes the bounded sine and at zero because $\sin 0=0$. Solving, $F(t)=F(0)\eu^{-t^{2}/4}$, and $F(0)=\tfrac{\sqrt{\pi}}{2}$ is the Gaussian integral, so
\[\int_0^{\infty}\eu^{-x^{2}}\cos(tx)\dd x=\frac{\sqrt{\pi}}{2}\,\eu^{-t^{2}/4}.\]
The formula is a machine. Since $\cos^{3}\theta=\tfrac34\cos\theta+\tfrac14\cos3\theta$ gives $\cos^{3}(2x)=\tfrac34\cos2x+\tfrac14\cos6x$, reading off $t=2$ and $t=6$ yields
\[\int_0^{\infty}\eu^{-x^{2}}\cos^{3}(2x)\dd x=\frac{\sqrt{\pi}}{2}\left(\frac34\eu^{-1}+\frac14\eu^{-9}\right)\approx0.244546,\]
confirmed to twelve places by quadrature. Any finite trigonometric polynomial in $\cos$ can be integrated against $\eu^{-x^{2}}$ this way, and the only work is the product to sum reduction.

A second instance, where the ODE comes from a substitution rather than from parts. Let $G(t)=\int_0^{\infty}\eu^{-x^{2}-t^{2}/x^{2}}\dd x$ for $t\geq0$. Differentiating,
\[G'(t)=-2t\int_0^{\infty}\frac{1}{x^{2}}\eu^{-x^{2}-t^{2}/x^{2}}\dd x,\]
and the substitution $u=t/x$, which sends $\dd x/x^{2}$ to $\dd u/t$ and leaves the exponent invariant, turns the right side into $-2G(t)$. So $G'=-2G$, and $G(0)=\tfrac{\sqrt\pi}{2}$ gives
\[\int_0^{\infty}\eu^{-x^{2}-t^{2}/x^{2}}\dd x=\frac{\sqrt{\pi}}{2}\,\eu^{-2t},\]
confirmed at $t=\tfrac12$ against the numerical value $0.326025$. The self reciprocal substitution $x\mapsto t/x$ is what made the ODE appear, which is a general lesson: an integral invariant under a substitution up to a factor will satisfy a differential equation in the parameter of that substitution.
\end{worked}

\begin{trap}
Discarding the boundary term without looking at it. Here $\bigl[\tfrac12\eu^{-x^{2}}\sin(tx)\bigr]_0^\infty$ vanishes, but in the same computation with $\eu^{-x}$ in place of $\eu^{-x^{2}}$ and a growing test function it does not, and the resulting ODE is simply false. The second trap is solving $I'=\lambda I$ and forgetting that a nonhomogeneous term $\mu(a)$ requires an integrating factor rather than a bare exponential.
\end{trap}

\begin{seealsob}Differentiation under the integral sign; Gaussian and error function integrals; cyclic integration by parts.\end{seealsob}

The ODE route is worth trying whenever the parameter sits in a place where $\partial_a$ multiplies the integrand by $x$ and where a factor of the integrand is the derivative of something simple. That is a narrow condition, but the family it covers is large: shifting $\cos(tx)$ to $\cosh(tx)$ turns the sign around and gives $\int_0^\infty \eu^{-x^{2}}\cosh(tx)\dd x=\tfrac{\sqrt\pi}{2}\eu^{t^{2}/4}$, and replacing $\eu^{-x^{2}}$ by $\eu^{-x^{2}/2}$ rescales the result rather than changing its shape. When the ODE is second order rather than first, you are in the theory of special functions and the parameter method has handed the problem to Bessel or to Airy, which is usually the right outcome.

\section{Frullani: two scales, one function}

The next technique looks at first like a different subject, but it is the same idea with the parameter differentiation already performed once and for all. A Frullani integrand compares one function at two different scales. Each half diverges logarithmically at one end, the difference converges, and the answer depends only on the two limiting values of $f$ and on the ratio of the scales. Nothing about the shape of $f$ in between matters at all, which is the surprising part and the reason the formula is worth memorising.

\tech{Frullani's integral}{medium}
\cue $\dfrac{f(ax)-f(bx)}{x}$ on $(0,\infty)$: the same function at two scales, over $x$. The signature is that each term alone would give a divergent $\int\frac{\dd x}{x}$ near $0$ or near $\infty$, and only the difference is finite.
\method If $f$ is continuous on $[0,\infty)$ and both $f(0^{+})$ and $f(\infty)$ exist and are finite, then for $a,b>0$
\[\int_0^{\infty}\frac{f(ax)-f(bx)}{x}\dd x=\bigl[f(0^{+})-f(\infty)\bigr]\ln\frac{b}{a}.\]
If instead $f(\infty)$ does not exist but $\int_A^\infty\frac{f(x)}{x}\dd x$ converges for some $A>0$, the same formula holds with $f(\infty)$ replaced by $0$. Those are the honest hypotheses; the formula is false without one of them.

\begin{worked}
The exponential case is the one to know by heart. With $f(x)=\eu^{-x}$, $f(0)=1$ and $f(\infty)=0$, so
\[\int_0^{\infty}\frac{\eu^{-x}-\eu^{-3x}}{x}\dd x=\ln3\approx1.098612.\]
It is also a parameter computation in disguise. Set $I(a)=\int_0^\infty\frac{\eu^{-ax}-\eu^{-3x}}{x}\dd x$; then $I'(a)=-\int_0^\infty \eu^{-ax}\dd x=-\tfrac1a$, so $I(a)=-\ln a+C$, and the anchor $I(3)=0$ gives $C=\ln3$ and $I(a)=\ln(3/a)$.

Two more, chosen to exercise the two hypothesis clauses. With $f=\arctan$, $f(0)=0$ and $f(\infty)=\tfrac{\pi}{2}$, so
\[\int_0^{\infty}\frac{\arctan(\pi x)-\arctan x}{x}\dd x=\Bigl(0-\frac{\pi}{2}\Bigr)\ln\frac{1}{\pi}=\frac{\pi}{2}\ln\pi\approx1.798137.\]
With $f=\cos$ there is no limit at infinity, but $\int_A^\infty\frac{\cos x}{x}\dd x$ converges by Dirichlet's test, so the second clause applies and $\int_0^\infty\frac{\cos(ax)-\cos(bx)}{x}\dd x=\ln\frac{b}{a}$; at $a=1$, $b=3$ this is $\ln3$ again, and the numerical partial integrals do converge to $1.09861$ once the oscillatory tail is averaged.
\end{worked}

\begin{keynote}[Why only the endpoints matter]
Cut the integral at $\varepsilon$ and $R$ and substitute $u=ax$ in the first term and $u=bx$ in the second. Both become $\int f(u)\frac{\dd u}{u}$ over shifted ranges, and the difference is
\[\int_{a\varepsilon}^{b\varepsilon}\frac{f(u)}{u}\dd u-\int_{aR}^{bR}\frac{f(u)}{u}\dd u.\]
On the first range $f$ is close to $f(0^{+})$ and $\int\frac{\dd u}{u}=\ln\frac{b}{a}$; on the second $f$ is close to $f(\infty)$ and the same logarithm appears. Everything between $\varepsilon$ and $R$ has cancelled exactly. That is the whole proof, and it explains both the form of the answer and why the hypotheses are about limits only.
\end{keynote}

Read from right to left, Frullani is a representation rather than an evaluation, and in that direction it manufactures $\gamma$. Setting $f=\eu^{-\cdot}$ and rearranging gives $\ln x=\int_0^{\infty}\frac{\eu^{-t}-\eu^{-xt}}{t}\dd t$, and multiplying by $\eu^{-x}$ and integrating in $x$ produces
\[\int_0^{\infty}\eu^{-x}\ln x\dd x=\int_0^{\infty}\frac{1}{t}\left(\eu^{-t}-\frac{1}{1+t}\right)\dd t=-\gamma\approx-0.577216,\]
the inner $x$ integral having been done first, which is legitimate here because the double integrand is absolutely integrable. The same object is $\Gamma'(1)=\psi(1)$, so the Euler constant enters the subject as the derivative of the Gamma function at $1$ and as a Frullani integral at the same time. That coincidence is why $\gamma$ turns up in parameter differentiation of Beta and Gamma integrals and almost nowhere else.

\subsection{The hypotheses, tested}

The formula has two hypothesis clauses and it is worth seeing all three cases that can arise. First, both limits exist: $f=\eu^{-\cdot}$, $f=\tfrac{1}{1+\cdot}$, $f=\arctan$, and the formula applies as stated. The middle case, $f(x)=\tfrac{1}{1+x}$, gives $\int_0^\infty\bigl(\tfrac{1}{1+x}-\tfrac{1}{1+5x}\bigr)\tfrac{\dd x}{x}=\ln5\approx1.609438$, which is a useful sanity check because the integrand is elementary and the antiderivative can be written down.

Second, $f(\infty)$ does not exist but $\int_A^\infty\frac{f}{x}$ converges. Both $\cos$ and $\sin$ are here. For $\cos$ the first clause value $f(0)=1$ survives and the answer is $\ln\frac{b}{a}$. For $\sin$ we have $f(0)=0$, so the answer is $0$, and indeed $\int_0^\infty\frac{\sin(ax)-\sin(bx)}{x}\dd x=\tfrac{\pi}{2}-\tfrac{\pi}{2}=0$ for any $a,b>0$: each half converges separately, which is the exceptional situation in this family and is why the $\sin$ case is a poor example to learn from.

Third, neither clause holds, and the integral diverges. This is the case that punishes formula matching, because the shape of the integrand is right and the answer produced is finite. The example in the trap below is the standard one.

\begin{trap}
Splitting the difference into two integrals. Both are infinite, so the split is not a manipulation of numbers at all; it is legal only inside a cut off $\int_\varepsilon^{R}$, where the divergences appear as $\ln$ terms that cancel. The other trap is applying the formula when neither hypothesis holds: for $f(x)=\ln(1+x)$ we have $f(0)=0$, so the naive answer is $0$, but $f(\infty)=\infty$ and $\int^\infty\frac{f}{x}$ diverges, and in fact $\int_0^\infty\frac{\ln(1+ax)-\ln(1+bx)}{x}\dd x$ diverges because the integrand behaves like $\frac{\ln(a/b)}{x}$ at infinity.
\end{trap}

\begin{seealsob}Choosing where to insert the parameter; epsilon regularisation of a divergent pair; Schwinger parametrisation.\end{seealsob}

\section{When the limits move}

So far the parameter has lived inside the integrand. It can also live in the limits, and then differentiation picks up boundary contributions. The general rule is Leibniz's, and its practical importance is less that it evaluates integrals directly than that it lets you build a function out of an integral, show its derivative is zero, and conclude that the integral is constant. That is a genuinely different mode of argument, and it produces the Gaussian integral out of nothing.

\tech{Leibniz rule with variable limits}{medium}
\cue Limits that depend on the differentiation variable, or a problem that asks for $\dvop{x}$ of an integral, or an integral appearing inside a larger expression whose derivative you want.
\method
\[\dvop{x}\int_{u(x)}^{v(x)}f(x,t)\dd t=f\bigl(x,v(x)\bigr)v'(x)-f\bigl(x,u(x)\bigr)u'(x)+\int_{u(x)}^{v(x)}\pdv{f}{x}\dd t.\]
All three terms are needed whenever $f$ carries the outer variable as well. Two of them are pure boundary bookkeeping; the third is differentiation under the sign and needs its own licence.

\begin{worked}
The Gaussian integral by a constancy argument. Define
\[\Phi(x)=\left(\int_0^{x}\eu^{-t^{2}}\dd t\right)^{2}+\int_0^{1}\frac{\eu^{-x^{2}(1+t^{2})}}{1+t^{2}}\dd t.\]
The first piece has a variable upper limit and no $x$ inside the integrand, so the chain rule and Leibniz give $2\eu^{-x^{2}}\int_0^x\eu^{-t^{2}}\dd t$. The second piece has fixed limits and $x$ inside, so only the third Leibniz term survives:
\[\dvop{x}\int_0^1\frac{\eu^{-x^{2}(1+t^{2})}}{1+t^{2}}\dd t=\int_0^1\frac{-2x(1+t^{2})\eu^{-x^{2}(1+t^{2})}}{1+t^{2}}\dd t=-2x\eu^{-x^{2}}\int_0^1 \eu^{-x^{2}t^{2}}\dd t,\]
and substituting $u=xt$ turns the last integral into $\tfrac1x\int_0^x \eu^{-u^{2}}\dd u$. The two derivatives cancel exactly, so $\Phi'\equiv0$ and $\Phi$ is constant. Evaluate at $x=0$: the first term is $0$ and the second is $\int_0^1\frac{\dd t}{1+t^{2}}=\tfrac{\pi}{4}$. Let $x\to\infty$: the second term is dominated by $\eu^{-x^{2}}$ and vanishes. Hence
\[\left(\int_0^{\infty}\eu^{-t^{2}}\dd t\right)^{2}=\frac{\pi}{4},\qquad \int_0^{\infty}\eu^{-t^{2}}\dd t=\frac{\sqrt{\pi}}{2}\approx0.886227.\]
No antiderivative was found, and no polar coordinates were used.
\end{worked}

The second standard use is structural rather than numerical: Leibniz turns a repeated integral into a single one. If $f$ is continuous, then
\[\dvop{x}\int_0^{x}(x-t)f(t)\dd t=\underbrace{(x-x)f(x)\cdot1}_{\text{boundary}}+\int_0^{x}f(t)\dd t=\int_0^{x}f(t)\dd t,\]
the boundary term vanishing because the integrand itself vanishes at $t=x$. Iterating gives Cauchy's formula for repeated integration,
\[\underbrace{\int_0^{x}\!\!\int_0^{t_1}\!\!\cdots\int_0^{t_{n-1}}}_{n\ \text{times}}f(t_n)\dd t_n\cdots\dd t_1=\frac{1}{(n-1)!}\int_0^{x}(x-t)^{n-1}f(t)\dd t,\]
which is the reason an $n$ fold antiderivative is a single convolution and the starting point for fractional integration. Both facts come out of the three term rule and nothing else.

\begin{trap}
Dropping the $\partial f/\partial x$ term because the limits were the interesting part. In the Gaussian argument that term is the entire content: without it $\Phi'$ is not zero and the argument collapses. A second and subtler failure is a moving limit that drags across a singularity of $f$, where the rule does not apply at all and the derivative can appear finite while the integral is not differentiable.
\end{trap}

\begin{seealsob}Differentiation under the integral sign; the parameter ODE; perfect derivative disguise.\end{seealsob}

The constancy argument generalises into a method. Whenever a problem gives you an integral together with a relation it satisfies, differentiate the relation: an integral equation of the form $y(x)=g(x)+\int_0^{x}K(x,t)y(t)\dd t$ becomes, after one application of the three term rule, the differential equation $y'(x)=g'(x)+K(x,x)y(x)+\int_0^{x}\partial_xK\,y\dd t$, and when $K$ does not depend on $x$ the integral term disappears entirely. That is the standard route from a Volterra equation to an initial value problem, and it is the same computation as the one that produced $\Phi'\equiv0$ above, run in the other direction.

\section{Parameters that live inside named functions}

The last two techniques exploit the fact that some parameters are already installed, with their derivatives already tabulated. The Beta function is a two parameter family whose derivatives are logarithms; the Gamma integral representation is a one parameter family that converts denominators into exponentials. Both are parameter differentiation, but the hard analytic work has been done once and recorded in the digamma function, so what remains is bookkeeping and a table of special values.

\tech{Beta and Gamma parameter differentiation}{hard}
\cue An integrand of Beta shape, $x^{p-1}(1-x)^{q-1}$ or $\sin^{2p-1}x\cos^{2q-1}x$, multiplied by $\ln x$, $\ln(1-x)$, $\ln\sin x$ or $\ln\cos x$. Logarithms of exactly those four kinds are what $\partial_p$ and $\partial_q$ produce, and nothing else does.
\method Start from
\[\Bfun(p,q)=\int_0^1t^{p-1}(1-t)^{q-1}\dd t=2\int_0^{\pi/2}\sin^{2p-1}x\,\cos^{2q-1}x\dd x=\frac{\Gamma(p)\Gamma(q)}{\Gamma(p+q)}.\]
Then $\partial_p$ inserts $\ln t$ in the algebraic form and $2\ln\sin x$ in the trigonometric one, and
\[\partial_p\partial_q\Bfun=\Bfun\Bigl\{\bigl[\psi(p)-\psi(p+q)\bigr]\bigl[\psi(q)-\psi(p+q)\bigr]-\psi'(p+q)\Bigr\}.\]

\begin{worked}
\[J=\int_0^{\pi/2}\frac{\ln\sin x\,\ln\cos x}{\sqrt{\tan x}}\dd x.\]
Write $\tan^{-1/2}x=\sin^{-1/2}x\cos^{1/2}x$, so $2p-1=-\tfrac12$ and $2q-1=\tfrac12$, giving $p=\tfrac14$, $q=\tfrac34$ and $p+q=1$. Since $\partial_p\partial_q$ applied to $2\int\sin^{2p-1}\cos^{2q-1}$ inserts $4\ln\sin x\ln\cos x$, we have $J=\tfrac18\,\partial_p\partial_q\Bfun$ at that point. Now $\Bfun(\tfrac14,\tfrac34)=\Gamma(\tfrac14)\Gamma(\tfrac34)=\frac{\pi}{\sin(\pi/4)}=\pi\sqrt2$ by Euler reflection, and with $\psi(1)=-\gamma$, $\psi(\tfrac14)=-\gamma-\tfrac{\pi}{2}-3\ln2$, $\psi(\tfrac34)=-\gamma+\tfrac{\pi}{2}-3\ln2$,
\[\bigl[\psi(\tfrac14)-\psi(1)\bigr]\bigl[\psi(\tfrac34)-\psi(1)\bigr]=\bigl(-\tfrac{\pi}{2}-3\ln2\bigr)\bigl(\tfrac{\pi}{2}-3\ln2\bigr)=9\ln^{2}2-\frac{\pi^{2}}{4}.\]
Subtracting $\psi'(1)=\tfrac{\pi^{2}}{6}$ and multiplying by $\Bfun/8$,
\[J=\frac{\pi\sqrt2}{8}\left(9\ln^{2}2-\frac{5\pi^{2}}{12}\right)\approx0.117593,\]
which quadrature confirms to fifteen places.
\end{worked}

The algebraic form of the Beta function is used the same way and is often simpler, because the logarithms it produces are $\ln t$ and $\ln(1-t)$ rather than $\ln\sin$ and $\ln\cos$. A single derivative gives
\[\int_0^1t^{p-1}(1-t)^{q-1}\ln t\dd t=\Bfun(p,q)\bigl[\psi(p)-\psi(p+q)\bigr],\]
so, for instance, $\int_0^1\frac{\ln t}{\sqrt{t(1-t)}}\dd t=\Bfun(\tfrac12,\tfrac12)\bigl[\psi(\tfrac12)-\psi(1)\bigr]=\pi(-2\ln2)=-2\pi\ln2\approx-4.355172$, which quadrature confirms. The corresponding second derivative in the same variable inserts $\ln^{2}t$ and brings in $\psi'$, and the mixed derivative is the formula quoted above. Three derivatives is where the method stops being pleasant and the polygamma bookkeeping starts to dominate; past that point, expand instead.

\begin{keynote}[The special values you need]
Parameter differentiation of $\Bfun$ always lands on $\psi$ at the arguments you started from, so the table is short. $\psi(1)=-\gamma$, $\psi(\tfrac12)=-\gamma-2\ln2$, $\psi(\tfrac14)=-\gamma-\tfrac{\pi}{2}-3\ln2$, $\psi(\tfrac34)=-\gamma+\tfrac{\pi}{2}-3\ln2$, and $\psi(n)=-\gamma+H_{n-1}$. For the trigamma, $\psi'(1)=\tfrac{\pi^{2}}{6}$, $\psi'(\tfrac12)=\tfrac{\pi^{2}}{2}$. Two structural facts generate the rest: the reflection $\psi(1-z)-\psi(z)=\pi\cot\pi z$, and the recurrence $\psi(z+1)=\psi(z)+\tfrac1z$. Notice that $\gamma$ cancels out of every difference $\psi(p)-\psi(p+q)$, which is why final answers contain $\pi$ and $\ln2$ but not $\gamma$.
\end{keynote}

\begin{trap}
The factor of $2$ in the trigonometric Beta form propagates into every derivative, so a double derivative carries $4$ and the prefactor is $\tfrac18$, not $\tfrac12$. The other reliable slip is sign: $\psi(\tfrac14)-\psi(1)$ is negative and $\psi(\tfrac34)-\psi(1)$ is negative too, since $3\ln2>\tfrac{\pi}{2}$, so their product is positive and the $-\psi'(p+q)$ term is what makes the bracket small.
\end{trap}

\begin{seealsob}The Beta function; digamma and polygamma values; differentiation under the integral sign.\end{seealsob}

The final technique runs the machinery backwards. Instead of differentiating with respect to a parameter, you introduce an integral over one, chosen so that the new double integral factorises after the order is swapped. It costs you a dimension and buys you a product structure, which is a good trade whenever the obstruction is a denominator that refuses to split. It is also the bridge to the next chapter, where swapping orders becomes the main event.

\tech{Schwinger parametrisation}{hard}
\cue A product of denominators, a denominator raised to a non-integer power, or a $\ln x$ or $1/x$ factor blocking a transform you would otherwise recognise.
\method Replace the obstruction by an integral and exchange the order. The four representations worth carrying are
\[\frac{1}{A^{s}}=\frac{1}{\Gamma(s)}\int_0^{\infty}t^{s-1}\eu^{-At}\dd t,\qquad \frac{1}{AB}=\int_0^1\frac{\dd t}{\bigl[tA+(1-t)B\bigr]^{2}},\]
\[\frac1x=\int_0^{\infty}\eu^{-xt}\dd t,\qquad \ln x=\int_0^{\infty}\frac{\eu^{-t}-\eu^{-xt}}{t}\dd t,\]
the last of which is Frullani read from right to left.

\begin{worked}
\[\int_0^{\infty}\frac{\dd x}{(1+x^{2})^{2}}.\]
Apply the first representation with $A=1+x^{2}$ and $s=2$, so $\frac{1}{(1+x^{2})^{2}}=\int_0^\infty t\,\eu^{-(1+x^{2})t}\dd t$. The double integrand is positive, so Tonelli licenses the swap unconditionally:
\[\int_0^{\infty}\!\!\int_0^{\infty}t\,\eu^{-t}\eu^{-tx^{2}}\dd t\dd x=\int_0^{\infty}t\,\eu^{-t}\left(\frac12\sqrt{\frac{\pi}{t}}\right)\dd t=\frac{\sqrt{\pi}}{2}\int_0^\infty t^{1/2}\eu^{-t}\dd t=\frac{\sqrt{\pi}}{2}\,\Gamma\!\left(\tfrac32\right)=\frac{\pi}{4}.\]
The point is structural: the representation turned a rational function of $x$ into a Gaussian in $x$ and a Gamma in $t$, and each factor was then standard.

The second representation handles products of unrelated denominators. With $A=x^{2}+a^{2}$ and $B=x^{2}+b^{2}$, the combined denominator is $x^{2}+M^{2}$ with $M^{2}=ta^{2}+(1-t)b^{2}$, and $\int_{-\infty}^{\infty}\frac{\dd x}{(x^{2}+M^{2})^{2}}=\frac{\pi}{2M^{3}}$. Hence
\[\int_{-\infty}^{\infty}\frac{\dd x}{(x^{2}+a^{2})(x^{2}+b^{2})}=\frac{\pi}{2}\int_0^1\frac{\dd t}{\bigl[ta^{2}+(1-t)b^{2}\bigr]^{3/2}}=\frac{\pi}{2(a^{2}-b^{2})}\left[\frac{2}{b}-\frac{2}{a}\right]=\frac{\pi}{ab(a+b)},\]
and at $a=2$, $b=3$ this is $\pi/30\approx0.104720$, which quadrature confirms. Partial fractions would also work here; the parametrisation is what survives when the powers are not integers.

The same representation proves a formula that would otherwise need a substitution and a Beta table. For $0<s<c$,
\[\int_0^{\infty}\frac{x^{s-1}}{(1+x)^{c}}\dd x=\frac{1}{\Gamma(c)}\int_0^{\infty}\!\!\int_0^{\infty}x^{s-1}t^{c-1}\eu^{-t}\eu^{-xt}\dd t\dd x=\frac{\Gamma(s)}{\Gamma(c)}\int_0^{\infty}t^{c-s-1}\eu^{-t}\dd t=\Bfun(s,c-s),\]
using $\int_0^\infty x^{s-1}\eu^{-xt}\dd x=\Gamma(s)t^{-s}$ in the middle step, and Tonelli throughout because everything is positive. At $s=\tfrac12$, $c=2$ this gives $\Bfun(\tfrac12,\tfrac32)=\tfrac{\pi}{2}$, and at $s=0.7$, $c=2.5$ it gives $0.909464$, both confirmed numerically.
\end{worked}

\begin{trap}
Swapping the order because it is convenient. Tonelli applies only to a nonnegative integrand; for a signed one you owe an absolute integrability check, and if the double integral is only conditionally convergent the two orders can differ. The second trap is using $\frac{1}{A^{s}}=\frac{1}{\Gamma(s)}\int t^{s-1}\eu^{-At}\dd t$ when $A$ is not positive, where the $t$ integral diverges outright.
\end{trap}

\begin{keynote}[Which representation to reach for]
Match the representation to the obstruction, not to the integrand. A denominator raised to an awkward power wants the Gamma representation, because it converts the power into a factor $t^{s-1}$ that the outer integral can absorb. A product of two unrelated denominators wants the Feynman representation, because it fuses them into one denominator at the cost of a single extra variable. A bare $1/x$ or $\ln x$ wants the exponential representations, because those trade a singular factor for an exponential that most integrands already like. In every case the test is the same: after the swap, is the inner integral a standard one? If not, you have added a dimension for nothing.
\end{keynote}

\begin{seealsob}Frullani's integral; justifying the swap; double integral and double series.\end{seealsob}

\subsection{When the family closes on itself}

Earlier, the family $\int_0^1\frac{\arctan(ax)}{x}\dd x$ reproduced itself and the protocol returned nothing. Self reference is not always fatal; sometimes it is the solution, because a linear equation in the unknown is as good as a formula. Take
\[I(a)=\int_0^1\frac{\ln(1+ax)}{1+x^{2}}\dd x,\qquad\text{target }I(1).\]
Differentiating and decomposing $\frac{x}{(1+ax)(1+x^{2})}=\frac{1}{1+a^{2}}\Bigl[\frac{-a}{1+ax}+\frac{x+a}{1+x^{2}}\Bigr]$,
\[I'(a)=\frac{1}{1+a^{2}}\left[-\ln(1+a)+\frac{\ln2}{2}+\frac{\pi a}{4}\right].\]
Now integrate from $0$ to $1$, using $I(0)=0$ and the two elementary values $\int_0^1\frac{\dd a}{1+a^{2}}=\frac{\pi}{4}$ and $\int_0^1\frac{a\dd a}{1+a^{2}}=\frac{\ln2}{2}$:
\[I(1)=-\int_0^1\frac{\ln(1+a)}{1+a^{2}}\dd a+\frac{\ln2}{2}\cdot\frac{\pi}{4}+\frac{\pi}{4}\cdot\frac{\ln2}{2}.\]
The first term on the right is $-I(1)$, since the integration variable is a dummy. Hence $2I(1)=\frac{\pi\ln2}{4}$ and
\[\int_0^1\frac{\ln(1+x)}{1+x^{2}}\dd x=\frac{\pi\ln2}{8}\approx0.272198,\]
confirmed numerically. The difference between this and the Catalan failure is exactly the sign with which the family reappears: a coefficient of $-1$ gives a solvable equation, a coefficient of $+1$ gives the tautology $I=I$. Always finish the calculation far enough to see which one you have.

The seven techniques of this chapter are one technique seen from seven angles. In each, an integral is embedded in a family, and the family is easier because it has a derivative. Feynman differentiates the family; the parameter ODE differentiates it and recognises the result as the family again; Frullani is the special case in which the differentiation has already been done in general; Leibniz extends the derivative to moving limits; Beta differentiation uses a family whose derivatives are tabulated; and Schwinger builds the family by integration instead of finding it by inspection. What follows in the next chapter is the same idea with a sum in place of the parameter integral.

Two limits of the method are worth stating plainly. Parameter differentiation cannot help when the integral is invariant along every direction you can insert a parameter, as the Catalan example above showed; the symptom is an identity rather than a value, and the cure is a different chapter. And it cannot help when $I'(a)$ has no closed form, which is the usual situation for integrands built from more than one transcendental factor. In that case the honest move is to stop after step 3 and read $I'(a)$ as a definition of a new function, which is how the polylogarithms and the elliptic integrals were discovered in the first place.

\begin{keynote}[Confirm before you commit]
Every closed form in this chapter is a two line numerical check: evaluate the family at two or three values of the parameter, compare against the formula, and only then believe it. The characteristic error in parameter work is not a wrong method but a lost constant of integration, and a lost constant shows up as an offset that is the same at every $a$. If your formula is wrong by a quantity that does not vary with $a$, the mistake is at step 6; if it is wrong by a quantity that does vary, the mistake is at step 3. Two evaluations distinguish the two cases in less time than rereading the derivation.
\end{keynote}

\section{Recognition matrix}
\begin{recog}{@{}p{0.30\textwidth}p{0.36\textwidth}p{0.28\textwidth}@{}}
\rowcolor{mist}\mh{If you see} & \mh{Reach for} & \mh{If that fails} \\
\midrule
\endhead
A constant that could be a letter & Differentiate under the sign; anchor where it is trivial & Insert a second parameter \\
$\dfrac{x^{a}-x^{b}}{\ln x}$ or any $\ln x$ downstairs & Parameter in the exponent, $\partial_a x^{a}=x^{a}\ln x$ & $\ln x=\int_0^\infty\frac{\eu^{-t}-\eu^{-xt}}{t}\dd t$ \\
$\dfrac{\sin x}{x}$, $\dfrac{1-\cos x}{x^{2}}$ on $(0,\infty)$ & Damping factor $\eu^{-ax}$, then let $a\to0^{+}$ & Justify $a\to0^{+}$ separately \\
$\arctan(ax)$ or $\ln(1+ax)$ over $x$ & Differentiate in $a$; the result is rational & Frullani, if it is a difference \\
$I'(a)$ reproduces $I(a)$ after parts & Solve the first order ODE in $a$ & Check the boundary term first \\
$\eu^{-x^{2}}\cos(tx)$ or $\eu^{-x^{2}}\cosh(tx)$ & $F'=\mp\tfrac t2F$, anchor at $F(0)=\tfrac{\sqrt\pi}{2}$ & Complete the square in the exponent \\
$\dfrac{f(ax)-f(bx)}{x}$ on $(0,\infty)$ & Frullani, $[f(0)-f(\infty)]\ln\frac{b}{a}$ & Check the second hypothesis, $\int^\infty f/x$ \\
A pair that each diverge logarithmically & Cut off at $\varepsilon$ and $R$, cancel, then take limits & Frullani with the Cesaro reading \\
Limits depending on $x$ & Leibniz with all three terms & Build $\Phi$ and show $\Phi'\equiv0$ \\
An $n$ fold repeated integral & Cauchy: $\frac{1}{(n-1)!}\int_0^x(x-t)^{n-1}f$ & Differentiate to check \\
$x^{p-1}(1-x)^{q-1}$ times a logarithm & $\partial_p$ or $\partial_q$ of $\Bfun(p,q)$ & Digamma values at quarter integers \\
$\sin^{m}x\cos^{n}x$ times $\ln\sin x$ & Trigonometric Beta; remember the $\tfrac18$ & Reflection $x\mapsto\tfrac{\pi}{2}-x$ \\
A product of two denominators & $\frac{1}{AB}=\int_0^1\frac{\dd t}{[tA+(1-t)B]^{2}}$ & Partial fractions in $x^{2}$ \\
$(1+x^{2})^{-s}$ with $s$ not an integer & $\frac{1}{A^{s}}=\frac{1}{\Gamma(s)}\int t^{s-1}\eu^{-At}\dd t$ & Substitute $x=\tan\theta$, then Beta \\
An answer that is off by a constant & Recheck the anchor, step 6 & Recompute $I'$ at one numeric $a$ \\
\end{recog}
